% shared by main.tex (the report) and ieee.tex (the IEEE version); edit here
A 180-degree U-bend in a differential pair makes the inner trace shorter by $\pi s$, and the usual fix is a
serpentine that adds the same length back. We test that fix with full-wave simulation instead of assuming it. We solve a controlled four-port fixture ($s = 0.508$\,mm, 20\,mil) in openEMS (FDTD) for \GeomCount{} geometries in
microstrip and stripline, \RunCount{} single-port-excited runs in all. Every result is taken through mixed-mode
S-parameters to receiver waveforms. The untuned bend gives \SkewB\,ps of skew, against \SkewSingle\,ps from geometry and the
single-line effective permittivity. Tuning that makes the two centrelines exactly equal in length removes only
\RemovedPctD--\RemovedPctF\,\% of that skew. \SkewF--\SkewD\,ps remain, and the shortfall grows with the number
of excursions. Where a tuning was solved on three or more meshes its value carries a mesh-based bound
($\pm$\UncC\,ps for C, $\pm$\UncD\,ps for D). The broad tuning F was solved on one mesh only. A transmission-line model built from the same lengths predicts
less than 1\,ps. In a serpentine experiment, folds closer together delivered less of
their predicted delay (delay efficiency \EtaSloose{} with folds 1.0\,mm apart, \EtaStight{} at 0.15\,mm). This is consistent with coupling between the folds. The shortfall that remains with loose folds is consistent with an
effect of the corners. Neither mechanism is isolated here. The same copper in stripline, solved on one mesh,
shows the same pattern. At 16\,Gb/s the
skew hardly closes the differential eye, from \EyeHA{} to \EyeHB\,\% of the swing, but turns \EyeVcB\,\% of the
swing into common mode. Sizing the tuning from its own residual in one more EM iteration brings the skew to \SkewFcal\ and
\SkewCcal\,ps for the broad and sparse tunings, but only to \SkewDcal\,ps for the dense accordion, consistent with its added amplitude coupling across its
folds too. Every result is a simulation from one solver, openEMS. Nothing was measured. A coupon plan for
measuring the predictions is given.
